\documentclass[portuguese, a4paper, 12pt, draft]{article}

\usepackage[brazil]{babel}
\usepackage[utf8]{inputenc}
\usepackage{amsmath,amsfonts}
\usepackage{newtxtext,newtxmath}
\usepackage{enumerate}
\usepackage{tabularx,multirow}
\usepackage{graphicx}
\usepackage[pdftex]{color}
\usepackage[pdftex,top=1.5cm, bottom=1.5cm, left=1.5cm, right=1.5cm]{geometry}
\usepackage{ifdraft}

% Using ABNT-like references
\usepackage[alf]{abntex2cite}

\newcommand{\R}{\mathbb{R}}

\begin{document}

\begin{figure}[h!]
  \centering
  \begin{tabular}{|cclc|}
    \hline
    \multirow{3}{*}{\includegraphics[width=0.15\textwidth,draft=false]{cnpq.png}} &
    \multirow{3}{*}{\includegraphics[width=0.15\textwidth,draft=false]{araucaria.png}} &
    \textcolor{blue}{\bf UNIVERSIDADE ESTADUAL DE MARINGÁ} \qquad \qquad &
    \multirow{3}{*}{\includegraphics[width=0.06\textwidth,draft=false]{uem.jpeg}} \\
    & & \textcolor{blue}{Pró-Reitoria de Pesquisa e Pós-Graduação} & \\
    & & \textcolor{blue}{\it Divisão de Pesquisa} & \\ \hline
  \end{tabular}
\end{figure}

\vspace{0.5cm}

\begin{center}
  \fbox{\makebox[0.98\textwidth]{\textcolor{blue}{\large \textbf{\textit{
            PROJETO DE INICIAÇÃO CIENTÍFICA -- PIBIC/CNPq-FA-UEM
          }}}}
  }
\end{center}

\vspace{1cm}

\ifdraft{
{\color{red}
Instruções: 
\begin{itemize}
    \item preencha os tópicos  abaixo com as informações pertinentes; 
    \item Após o preenchimento deste formulário gere uma versão em PDF deste arquivo, que deverá ser anexada ao SGP – Sistema de Gestão de Projetos, no local indicado durante o cadastro do projeto de iniciação científica - PIBIC no sítio http://www.sgp.uem.br. Deverão ser anexadas, também, as versões em PDF dos currículos Lattes, atualizados,  do orientador e  do aluno candidato a bolsista;
    \item Utilize Fonte Times New Roman 12, espaço simples e parágrafo justificado (já ajustado neste modelo).
    \item Ao finalizar, remova a opção \texttt{draft} da primeira linha e substitua por \texttt{final}.
\end{itemize}

\textbf{Obs.}: o preenchimento deste formulário é apenas uma etapa dos procedimentos para submissão de Projetos PIBIC, sendo necessário cadastrar o projeto e outras informações  no sítio www.sgp.uem.br. }
}

\begin{center}
  \begin{tabularx}{\textwidth}{|X|X|}
    \hline
    \multicolumn{2}{|>{\hsize=\dimexpr2\hsize+2\tabcolsep+\arrayrulewidth\relax}X|}
    {\textbf{1. TÍTULO DO PROJETO:} Otimização sem derivadas aplicada a um rastreador solar} \\ \hline
    \textbf{2. INÍCIO}: \textcolor{blue}{\bf 01/08/2019} &
    \textbf{3. FIM}: \textcolor{blue}{\bf 31/07/2020} \\ \hline
    \multicolumn{2}{|>{\hsize=\dimexpr2\hsize+2\tabcolsep+\arrayrulewidth\relax}X|}
    {\textbf{4. EQUIPE EXECUTORA}:} \\ \hline
    \multicolumn{2}{|>{\hsize=\dimexpr2\hsize+2\tabcolsep+\arrayrulewidth\relax}X|} {a) Nome do acadêmico: Matheus de Sousa Rodrigues} \\ \hline
    \multicolumn{2}{|>{\hsize=\dimexpr2\hsize+2\tabcolsep+\arrayrulewidth\relax}X|}
    {b) Nome do orientador: Francisco Nogueira Calmon Sobral} \\
    \hline
    \multicolumn{2}{|>{\hsize=\dimexpr2\hsize+2\tabcolsep+\arrayrulewidth\relax}X|}
    {c) Nome do coorientador: Emerson Vitor Castelani} \\ \hline
  \end{tabularx}
\end{center}

\setcounter{section}{4}

\ifdraft{\section{Resumo \textcolor{red}{(até 3000 caracteres)}}}{\section{Resumo}}

\noindent
\textbf{Palavras-chave:} 

\section{Introdução}


\section{Justificativas}

\section{Objetivos}

\section{Metodologia}


\section{Plano de trabalho individual}

\section{Cronograma de execução}

\begin{center}
  \begin{tabularx}{0.98\textwidth}{|X||c|c|c|c|c|c|c|c|c|c|c|c|}
\hline
\multicolumn{13}{|c|}{\textsc{Cronograma de execução}} \\ \hline
Descrição das atividades &
\multicolumn{12}{|c|}{Assinalar o mês em que a atividade será executada} \\ \hline
 & 1$^\circ$ & 2$^\circ$ & 3$^\circ$ & 4$^\circ$ & 5$^\circ$ & 6$^\circ$ &
 7$^\circ$ & 8$^\circ$ & 9$^\circ$ & 10$^\circ$ & 11$^\circ$ & 12$^\circ$ \\ \hline
       &X&X&X& & & & & & & & & \\ \hline
       & & &X&X&X& & & & & & & \\ \hline
       & & & & & &X&X& & & & & \\ \hline
       & & & &X&X&X&X&X&X&X&X& \\ \hline
       & & & & & &X& & & & &X&X\\ \hline
\end{tabularx}
\end{center}

\bibliography{pibic}
\bibliographystyle{modabnt}


\end{document}
